\subsection{COUNTABLE SETS}

DEFINITION 3. A set is said to be countable if it is equipotent to a subset of the set N of integers.

PROPOSITION 1. Every subset of a countable set is countable. The product of a finite family of countable sets is countable. The union of a sequence of countable sets is countable.

The first assertion is obvious. The others follow from the Corollaries to Theorem 2, no. 3.

We have proved (no. 3, Lemma 1) that if \(\mathfrak{a}\) is any infinite cardinal, then \(\operatorname{Card}(\mathbf{N}) \leqslant \mathfrak{a}\) . This has the following consequences:

PROPOSITION 2. Every countable infinite set E is equipotent to N.

For \(\operatorname{Card}(\mathbf{E}) \leqslant \operatorname{Card}(\mathbf{N})\) by definition; and \(\operatorname{Card}(\mathbf{N}) \leqslant \operatorname{Card}(\mathbf{E})\) since \(\mathbf{E}\) is infinite.

PROPOSITION 3. Every infinite set has a partition \((\mathbf{X}_i)_{i\in I}\) formed of countable infinite sets \(\mathbf{X}_i\) , the index set I being equipotent to E.

For \(\operatorname{Card}(\mathbf{E}) = \operatorname{Card}(\mathbf{E})\) \(\operatorname{Card}(\mathbf{N})\) (no. 3, Theorem 2, Corollary 4).

PROPOSITION 4. Let \(f\) be a mapping of a set \(\mathbf{E}\) onto an infinite set \(\mathbf{F}\) such that, for each \(y \in \mathbf{F}\) , \(\overline{f}^1(y)\) is countable. Then \(\mathbf{F}\) is equipotent to \(\mathbf{E}\) .

For the sets \(\overline{f}^1 (y)\) \((y\in \mathbf{F})\) form a partition of \(\mathbf{E}\) ; hence

\[
\operatorname{Card} (\mathbf {E}) \leqslant \operatorname{Card} (\mathbf {F}) \operatorname{Card} (\mathbf {N}) = \operatorname{Card} (\mathbf {F}),
\]

and \(\operatorname{Card}(\mathbf{F}) \leqslant \operatorname{Card}(\mathbf{E})\) by Proposition 3 of §3, no. 2.

PROPOSITION 5. The set \(\mathfrak{F}(\mathbf{E})\) of finite subsets of an infinite set \(\mathbf{E}\) is equipotent to \(\mathbf{E}\) .

For each integer \(n\) , let \(\mathfrak{F}_n\) denote the set of subsets of \(\mathbf{E}\) which have \(n\) elements. For each \(\mathbf{X} \in \mathfrak{F}_n\) there exists a bijection of \([1, n]\) onto \(\mathbf{X}\) . Hence the cardinal of \(\mathfrak{F}_n\) is at most equal to that of the set of mappings of \([1, n]\) into \(\mathbf{E}\) , i.e., to \(\operatorname{Card}(\mathbf{E}^n) = \operatorname{Card}(\mathbf{E})\) (no. 3, Theorem 2, Corollary 1). Therefore

\[
\operatorname{Card} (\mathfrak {F} (E)) = \sum_ {n \in N} \operatorname{Card} (\mathfrak {F} _ {n}) \leqslant \operatorname{Card} (E) \operatorname{Card} (N) = \operatorname{Card} (E).
\]

On the other hand, since \(x \to \{x\}\) is an injective mapping of \(\mathbf{E}\) into \(\mathfrak{F}(\mathbf{E})\) , we have \(\operatorname{Card}(\mathbf{E}) \leqslant \operatorname{Card}(\mathfrak{F}(\mathbf{E}))\) .

DEFINITION 4. A set is said to have the power of the continuum if it is equipotent to the set of all subsets of N.

A set which has the power of the continuum is not countable (§ 3, no. 6, Theorem 2).

* The name ``power of the continuum'' arises from the fact that the set of real numbers is equipotent to \(\mathfrak{P}(\mathbf{N})\) (General Topology, Chapter IV, §8). * The continuum hypothesis is the assertion that every uncountable set contains a subset which has the power of the continuum; the generalized continuum hypothesis is the assertion that, for every infinite cardinal \(a\) , every cardinal \(>a\) is \(\geqslant 2^{a}\) .

